% !TeX root = Asparsi.tex
% !BIB TS-program = biber
% !TeX encoding = UTF-8
% !TeX spellcheck = it_IT
\chapter{Racconti}
\section{Aritmetica}
\begin{racconto}[De Giorgi]
Un giorno don Raffaele Ramirez propose il seguente problema alla sua classe di terza elementare della scuola "Cesare Battisti" in via Achille Costa, a Lecce.
"In una fattoria ci sono oche e conigli.
Le oche, come sapete, hanno 2 zampe, mentre i conigli ne hanno 4.
In tutto, nella fattoria ci sono 11 animali, che in totale hanno 34 zampe.
Quanti sono i conigli e quante sono le oche?".
Qualcuno tra i banchi ha la risposta pronta, esita per un istante e poi alza timidamente la mano.
"Che cosa c'è Ennio? Non è chiaro?", domanda il maestro.
"Ci sono 5 oche e 6 conigli", risponde il bambino.
"Bravo Ennio. 
Ma come hai fatto a rispondere così in fretta?
Conoscevi già la soluzione?".
"No", ribatte Ennio, "ma il problema è semplice".
"Davvero? E come si fa a risolverlo?", chiede perplesso il maestro.
"Se prendo 34 e divido per 2 ottengo 17, che è pari alla somma dei bipedi e dei quadrupedi contati due volte.
Ma so già che 11 è la somma dei bipedi e dei quadrupedi, quindi la differenza tra 17 e 11, cioè 6, è pari al numero dei quadrupedi.
E i bipedi devono essere 5".
Le premesse c'erano già tutte, e in effetti quel bambino, Ennio De Giorgi\index{De Giorgi!Ennio}, sarebbe diventato uno dei più geniali matematici italiani del XX secolo.
\end{racconto}
\[
\left \{\begin{array}{ll}
x+y=&11\\
2x+4y=&34
\end{array}
\right.
\]